% !TEX root = ../../ctfp-print.tex

\lettrine[lhang=0.17]{モ}{ノイドは圏論においても}プログラミングにおいても重要な概念だ。
圏は強い型付け言語に対応し、モノイドは型なし言語に対応する。
型なし言語ではどんな二つの関数も合成できるように、モノイドではどんな二つの射も合成できるからだ
（もちろん、プログラムを実行するときに実行時エラーが発生することもあるが）。

モノイドは単一の対象を持つ圏として記述でき、そこではすべてのロジックが射の合成規則に符号化されている
ことを見てきた。この圏論的モデルは、集合論的なモノイドの伝統的な定義と完全に等価だ。
集合論的定義では、集合の二つの元を「掛け合わせて」第三の元を得る。この「掛け算」の過程は、
まず元のペアを形成し、そのペアを既存の元――その「積」――と同一視するという二段階にさらに分解できる。

掛け算の第二部分――ペアを既存の元と同一視する操作――を省いたらどうなるだろうか？
例えば、任意の集合から始め、元のすべてのペアを形成し、それらを新しい元と呼ぶことができる。
次にこれらの新しい元をすべての可能な元とペアを作り、それをまた繰り返す。
これは連鎖反応だ――新しい元を永遠に追加し続けることになる。
結果として得られる無限集合は、\emph{ほぼ}モノイドだ。
しかしモノイドには単位元と結合律も必要だ。問題ない。特別な単位元を追加し、
単位律と結合律を支えるのに十分なだけペアのいくつかを同一視すればよい。

簡単な例でこれがどのように機能するか見てみよう。二つの元の集合 $\{a, b\}$ から始めよう。
これらを自由モノイドの\newterm{生成元}と呼ぶ。まず、単位元として機能する特別な元 $e$ を追加する。
次に元のすべてのペアを追加し、それらを「積」と呼ぶ。$a$ と $b$ の積はペア $(a, b)$ になる。
$b$ と $a$ の積はペア $(b, a)$、$a$ と $a$ の積は $(a, a)$、$b$ と $b$ の積は $(b, b)$ だ。
$e$ を含むペア、例えば $(a, e)$、$(e, b)$ なども形成できるが、これらはそれぞれ $a$、$b$ などと
同一視する。だからこのラウンドでは $(a, a)$、$(a, b)$、$(b, a)$、$(b, b)$ だけを追加し、
集合 $\{e, a, b, (a, a), (a, b), (b, a), (b, b)\}$ を得る。

\begin{figure}[H]
  \centering
  \includegraphics[width=0.8\textwidth]{images/bunnies.jpg}
\end{figure}

\noindent
次のラウンドでは $(a, (a, b))$、$((a, b), a)$ などの元を追加し続ける。
この時点で結合律が成り立つことを確認しなければならないので、
$(a, (b, a))$ と $((a, b), a)$ などを同一視する。言い換えれば、内側の括弧は不要になる。

このプロセスの最終結果は何か予想できるだろう。$a$ と $b$ の可能なすべてのリストを生成することになる。
実際、$e$ を空リストとして表すと、私たちの「掛け算」はリストの連結に他ならないことがわかる。

このように、元のあらゆる組み合わせを生成し続け、法則を支えるのに必要な最低限の同一視だけを行う構成を
\newterm{自由構成}と呼ぶ。今行ったのは、生成元の集合 $\{a, b\}$ から\newterm{自由モノイド}を構成することだ。

\section{Haskellにおける自由モノイド}

Haskellにおける二元集合は型 \code{Bool} と等価であり、この集合によって生成される自由モノイドは
型 \code{{[}Bool{]}}（\code{Bool} のリスト）と等価だ
（無限リストの問題は意図的に無視している）。

Haskellのモノイドは型クラスによって定義される：

\src{snippet01}
これは、すべての \code{Monoid} が \code{mempty} と呼ばれる単位元と、
\code{mappend} と呼ばれる二項関数（掛け算）を持たなければならないことを示している。
単位律と結合律はHaskellで表現できないため、モノイドをインスタンス化するたびにプログラマが
検証しなければならない。

任意の型のリストがモノイドを形成するという事実は、このインスタンス定義によって記述される：

\src{snippet02}
空リスト \code{{[}{]}} が単位元であり、リストの連結 \code{(++)} が二項演算であることを示している。

見てきたように、型 \code{a} のリストは、集合 \code{a} を生成元とする自由モノイドに対応する。
掛け算を持つ自然数の集合は自由モノイドではない。なぜなら多くの積を同一視しているからだ。
例えば比べてみよう：

\src{snippet03}
簡単だが、問題は、圏論ではオブジェクトの内部を見ることが許されていないのに、
この自由構成を圏論で行えるかどうかだ。私たちの汎用ツール：普遍的構成を使おう。

第二の興味深い問いは、法則が要求する最低限以上の元の同一視を行うことで、
任意のモノイドを何らかの自由モノイドから得られるかどうかだ。これが普遍的構成から直接従うことを示そう。

\section{自由モノイドの普遍的構成}

普遍的構成の以前の経験を思い出すと、それは何かを構成するというよりも、
与えられたパターンに最もよく適合するオブジェクトを選択することだと気づくかもしれない。
だから普遍的構成を使って自由モノイドを「構成」したいなら、選択肢となるモノイドの全体を
考慮しなければならない。選択できるモノイドの圏全体が必要だ。しかしモノイドは圏を形成するだろうか？

まず、モノイドを単位元と掛け算によって定義された追加構造を持つ集合として見てみよう。
射としては、モノイド構造を保存する関数を選ぶ。このような構造保存関数を\newterm{準同型写像}と呼ぶ。
モノイド準同型写像は、二つの元の積をその二つの元の写像の積に写さなければならない：

\src{snippet04}
また、単位元を単位元に写さなければならない。

例えば、整数のリストから整数へのモノイド準同型写像を考えよう。
\code{{[}2{]}} を 2 に、\code{{[}3{]}} を 3 に写すなら、連結

\src{snippet05}
が掛け算

\src{snippet06}
になるように \code{{[}2, 3{]}} を 6 に写さなければならない。

では、個々のモノイドの内部構造を忘れて、対応する射を持つ対象としてのみ見てみよう。
モノイドの圏 $\cat{Mon}$ が得られる。

内部構造を忘れる前に、重要な性質に気づいておこう。$\cat{Mon}$ のすべての対象は自明に集合に写せる。
それは単にその元の集合だ。この集合を\newterm{台集合}と呼ぶ。実際、$\cat{Mon}$ の対象を
集合に写せるだけでなく、$\cat{Mon}$ の射（準同型写像）も関数に写せる。
これも自明に思えるが、すぐに役立つ。$\cat{Mon}$ から $\Set$ への対象と射のこの写像は
実際には関手だ。この関手はモノイド構造を「忘れる」――普通の集合の内部に入ってしまうと、
単位元を区別したり掛け算を気にしたりしなくなる――ので、\newterm{忘却関手}と呼ばれる。
忘却関手は圏論に頻繁に登場する。

$\cat{Mon}$ についての二つの異なる見方ができるようになった。
対象と射を持つ他の圏と同様に扱うことができる。その見方では、モノイドの内部構造は見えない。
$\cat{Mon}$ の特定の対象について言えることは、恒等射と他の対象への射を通じて自分自身や
他の対象とつながっているということだけだ。射の「掛け算表」――合成規則――は
もう一方の見方：集合としてのモノイドから導出される。圏論に移行することでこの見方を
完全に失ったわけではない――忘却関手を通じてまだアクセスできる。

普遍的構成を適用するには、モノイドの圏を探索して自由モノイドの最良候補を選べるような
特別な性質を定義する必要がある。しかし自由モノイドはその生成元によって定義される。
生成元の異なる選択は異なる自由モノイドを生む（\code{Bool} のリストは \code{Int} のリストと同じではない）。
この構成は生成元の集合から始まらなければならない。つまり集合に戻ってきた！

そこで忘却関手が登場する。モノイドをX線で透過して見るのに使える。
そのX線像の中で生成元を特定できる。仕組みはこうだ：

生成元の集合 $x$ から始める。これは $\Set$ における集合だ。

照合するパターンは、モノイド $m$――$\cat{Mon}$ の対象――と $\Set$ における関数 $p$ から成る：

\src{snippet07}
ここで $U$ は $\cat{Mon}$ から $\Set$ への忘却関手だ。これは奇妙な不均一なパターンだ――
半分は $\cat{Mon}$、半分は $\Set$ にある。

アイデアは、関数 $p$ が $m$ のX線像の中で生成元の集合を特定するというものだ。
関数が集合内の点を同一視する（潰す）ことがあっても問題ない。
すべては普遍的構成によって整理され、このパターンの最良の代表者が選ばれる。

\begin{figure}[H]
  \centering
  \includegraphics[width=0.4\textwidth]{images/monoid-pattern.jpg}
\end{figure}

\noindent
候補間のランク付けも定義しなければならない。別の候補：モノイド $n$ と
そのX線像の中で生成元を特定する関数があるとしよう：

\src{snippet08}
$m$ は、モノイドの射（構造保存準同型写像）

\src{snippet09}
が存在して、その $U$ による像（$U$ は関手なので射を関数に写す）が $p$ を通じて因子分解されるとき、
$n$ より良いと言う：

\src{snippet10}
$p$ が $m$ の中の生成元を選択し、$q$ が $n$ の中の「同じ」生成元を選択すると考えると、
$h$ はこれらの生成元を二つのモノイド間で写すと考えられる。$h$ は定義によってモノイド構造を
保存することを忘れないように。これは、一方のモノイドにおける二つの生成元の積が、
他方のモノイドにおける対応する二つの生成元の積に写されることを意味する。

\begin{figure}[H]
  \centering
  \includegraphics[width=0.4\textwidth]{images/monoid-ranking.jpg}
\end{figure}

\noindent
このランク付けを使って最良の候補――自由モノイド――を見つけられる。定義はこうだ：

\begin{quote}
  $m$（関数 $p$ とともに）が生成元 $x$ を持つ\textbf{自由モノイド}であるのは、
  上記の因子分解性質を満たす $m$ から他の任意のモノイド $n$（関数 $q$ とともに）への
  \emph{一意な}射 $h$ が存在する場合に限る。
\end{quote}
ちなみに、これは第二の問いに答えている。関数 $U h$ は $U m$ の複数の元を $U n$ の
一つの元に潰す力を持つ。この潰しは自由モノイドの一部の元を同一視することに対応する。
したがって、生成元 $x$ を持つ任意のモノイドは、$x$ に基づく自由モノイドからいくつかの元を
同一視することで得られる。自由モノイドは最低限の同一視だけが行われたものだ。

自由モノイドについては随伴を話すときに戻ってくる。

\section{課題}

\begin{enumerate}
  \tightlist
  \item
        モノイドの準同型写像が単位元を保存するという要件は冗長だと思うかもしれない（私も最初はそう思った）。
        なぜなら、すべての $a$ について

        \begin{snip}{text}
h a * h e = h (a * e) = h a
\end{snip}
        が成り立つことがわかっているので、$h e$ は右単位元として（また類比から左単位元としても）
        振る舞うからだ。問題は、すべての $a$ についての $h a$ が対象モノイドの部分モノイドしか
        覆わないかもしれないことだ。$h$ の像の外に「真の」単位元が存在する可能性がある。
        掛け算を保存するモノイド間の同型写像が単位元を自動的に保存しなければならないことを示せ。
  \item
        連結を持つ整数のリストから、掛け算を持つ整数へのモノイド準同型写像を考えよう。
        空リスト \code{{[}{]}} の像は何か？すべての単元リストはそれが含む整数に写されると仮定する、
        すなわち \code{{[}3{]}} は 3 に写されるなど。\code{{[}1, 2, 3, 4{]}} の像は何か？
        整数 12 に写されるリストはいくつあるか？二つのモノイド間に他の準同型写像はあるか？
  \item
        一元集合によって生成される自由モノイドは何か？それが何と同型かわかるか？
\end{enumerate}
